% Academic paper: integrating robotics models and control stacks into one
% multi-agent testbed, with external baselines at each stage.
% The verification of the base system against C++ SCRIMMAGE is baseline.tex.
% Choose the venue and replace prompts after the experiments establish results.
\documentclass[conference]{IEEEtran}
\usepackage{times}
\usepackage[numbers]{natbib}
\usepackage{booktabs}
\usepackage{tabularx}
\usepackage{array}
% Ragged-right paragraph column for narrow table cells.
\newcolumntype{P}[1]{>{\raggedright\arraybackslash}p{#1}}
\usepackage[bookmarks=true]{hyperref}

\newcommand{\stub}[1]{\par\noindent\textit{TODO: #1}\par}
% Marks results that do not yet come from a release run of the suite
% (reference/suite.py).
\newcommand{\prov}{\textbf{[Provisional]}~}

\begin{document}

\title{Ripple: Tracing the Effects of Model Fidelity in Multi-Agent Robotics}
\author{Author names and affiliations to be added}
\maketitle

\begin{abstract}
Small differences in simulated vehicle behavior can propagate through sensing,
communication, and decisions to change a team's outcome. Determining when
these differences matter requires controlled comparisons across levels of
model fidelity. We present Ripple, a Rust research testbed for tracing these
effects, which brings fixed-wing aircraft, multirotors, VTOL aircraft, and
surface vessels into a shared multi-agent world; fixed-wing aircraft with
aerodynamic tables, terrain contact, and an external autopilot are built, and
the others are in progress.
Researchers select vehicle dynamics, aerodynamic and hydrodynamic effects,
and native or external autopilot control while retaining a common sensing,
communication, and estimation pipeline. Simple and detailed vehicles can
coexist in the same experiment. Each integration is checked against an
external implementation or published model, with explicit provenance for its
equations and parameters. The evaluation distinguishes agreement with a source
from physical validation against independent measurements and from the effect
of fidelity on a multi-agent conclusion. Controlled comparisons vary model
effects and control stacks, recording task outcomes, operating regimes, and
computational cost. \textit{[Results and findings pending: identify cases where
an apparently successful low-fidelity experiment develops catastrophic
failures or qualitatively different team behavior with added fidelity.
Quantify the change and trace it to the responsible model effects and
interactions.]}
The testbed provides a common experimental setting for measuring how local
model differences propagate, how much they change multi-agent results, and
when simpler models are sufficient.
\end{abstract}

\IEEEpeerreviewmaketitle

\section{Introduction}

Multi-agent robotics experiments connect decisions made by individual robots
to outcomes for a team. The vehicle and environment models determine how
those decisions become motion: actuator lag delays a response, aerodynamic
stall limits a manoeuvre, and waves move a vessel carrying a sensor. A small
change in motion can change when a robot detects a target, which report it
shares, and how another robot responds. Through these interactions, a local
difference can cascade into a different team outcome; it can also diminish or
leave the outcome effectively unchanged.

Studying these effects requires a setting in which researchers can change the
physical model while preserving the task, coordination strategy, and sensing
and communication rules. The motivating question is: \emph{when, through
which interactions, and by how much do differences in vehicle and environment
fidelity change the results of a multi-agent experiment?}

Answering this question requires substantial integration. Vehicle models,
autopilots, sensors, and communication models bring different conventions for
frames, units, actuator commands, and time. A component can reproduce its
source in isolation yet behave differently once scheduling, feedback, and
communication are introduced. A successful demonstration alone cannot show
whether the integrated component preserves the source behavior. Evidence is
needed both at the component boundary and in the assembled system.

We present Ripple, a research testbed that brings fixed-wing aircraft,
multirotors, VTOL aircraft, and surface vessels into one simulation with
selectable levels of model detail. Its Rust core builds on SCRIMMAGE's plugin
design \cite{demarco2018}, with a separate comparison establishing agreement
with the C++ baseline \cite{baselinereport}. Aerodynamic tables, terrain
contact, and ArduPilot in lockstep are integrated; rotor dynamics and vessel
response to waves are planned (Section~\ref{sec:platforms}). Native controllers
and external autopilots operate through explicit
state and command interfaces. Vehicles share sensing, communication, and
estimation stages that keep physical truth separate from each agent's belief.
This lets a few detailed vehicles operate alongside many simple ones without
changing the experiment's information rules.

Trust is established separately for each integration. Comparisons against
external implementations and published models check that equations, interface
conversions, and timing retain their intended behavior. Parameter provenance
and comparisons with independent measurements establish the extent of the
physical evidence. These are distinct claims: reproducing a reference model
does not make its estimated coefficients measured, and adding an autopilot
does not validate the plant it controls. Keeping these distinctions explicit
allows researchers to choose models according to the experiment and the
evidence supporting them.

The paper makes three contributions:
\begin{enumerate}
\item A common execution and information architecture that integrates
heterogeneous vehicles, selectable physical effects, and external control
stacks in repeatable multi-agent experiments.
\item External baseline comparisons for the integrations, with traceable
equations and parameters and explicit limits on what source agreement and
physical measurements establish.
\item A controlled evaluation of fidelity choices through task outcomes,
operating regimes, and computational cost, allowing both changed conclusions
and agreement between simpler and more detailed models to be assessed.
\end{enumerate}

The evaluation treats model structure, parameter evidence, and control stack
as separate experimental choices. It measures outcome differences and
strategy rankings alongside the physical regimes that produce them. Of
particular interest are cases in which a strategy appears successful under a
simple model but fails catastrophically or produces qualitatively different
team behavior when additional physical effects are included. Such cases can
show how a local modeling assumption changes observations and decisions until
the team reaches a different outcome.

\stub{Replace this placeholder with the strongest demonstrated cases and
their findings: the apparent success at low fidelity, the failure or changed
behavior with added fidelity, and the effects and interactions responsible.
The current landing study is preliminary evidence, not the headline example;
do not assume a catastrophic difference before the experiments establish it.}

Section~\ref{sec:system} describes the shared simulation architecture, and
Section~\ref{sec:verification} presents the integration comparisons.
Sections~\ref{sec:fidelity} and~\ref{sec:provenance} describe fidelity choices
and model provenance. Section~\ref{sec:stall} examines the fixed-wing case;
Sections~\ref{sec:platforms} and~\ref{sec:fleets} describe the vehicle models
and their composition into heterogeneous fleets. The coordination workloads
and policy interfaces follow in Sections~\ref{sec:workloads}
and~\ref{sec:learning}.

\section{Related Work}
\label{sec:related}

\paragraph{Simulators for many agents}
PettingZoo standardized how a learner steps a multi-agent environment
\cite{terry2021pettingzoo}. Isaac Gym, VMAS, and JaxMARL moved simulation onto
the GPU, advancing thousands of independent worlds together
\cite{makoviychuk2021isaacgym,bettini2022vmas,rutherford2023jaxmarl}; the
constraint is computation expressed as regular array operations, which MJX
meets by batching parameters across worlds \cite{mjxdocs}, and Madrona and
GPUDrive by an entity-component engine
\cite{shacklett2023madrona,kazemkhani2025gpudrive}.
Since 2025 these tools have moved toward physical fidelity: Isaac Lab adds
actuator models and multi-frequency sensors
\cite{mittal2025isaaclab,isaaclab3backends},
Crazyflow simulates drone swarms in JAX \cite{schuck2026crazyflow}, and
MuJoCo-Drones-Gym models aerodynamic effects on the CPU while its GPU backend
omits them \cite{tayal2026mujocodrones}. A parallel line keeps fast
environments on the CPU and only learning on the GPU
\cite{petrenko2020samplefactory,suarez2025pufferlib}. Communication is usually
an action channel or part of the observation; queued delivery with delay and
loss comes from co-simulation with a network simulator
\cite{burmeister2025defiance}. Multi-agent testbeds built for many vehicles,
such as SCRIMMAGE \cite{demarco2018}, run simple models and delegate physics to
a separate flight model where needed.

\paragraph{Vehicle dynamics and domain simulators}
JSBSim builds aircraft from coefficient tables \cite{berndt2004jsbsim}.
ArduPilot and PX4 ship their own simulators to test autopilot code
\cite{ardupilotsitl,px4gazebo}, and RotorS models multirotor motors and rotors
in Gazebo \cite{furrer2016rotors}. Fossen's Marine Systems Simulator gives
vessel manoeuvring and seakeeping models \cite{fossen2021handbook}; VRX
simulates an uncrewed surface vessel in waves and wind for an autonomy
competition \cite{bingham2019vrx}; panel solvers such as Capytaine compute
hydrodynamic coefficients from a hull \cite{ancellin2019capytaine}.

\paragraph{Fidelity and evaluation}
Robotics projects increasingly separate fast training from a more realistic
evaluation stage. Pyquaticus trains maritime game policies in a lightweight
simulator with embedded vehicle dynamics and deploys them on uncrewed surface
vessels \cite{crowley2025pyquaticus}. A connected-vehicle benchmark separates
degradation caused by differences between control stacks from degradation
caused by increasing environmental realism \cite{beerwerth2026zeroshot}, and
BVR Sim shows a policy transferring across aircraft through a common
controller-level interface \cite{sun2026bvrsim}. SMACv2 shows that success on a
fixed benchmark may not require the observations it provides
\cite{ellis2023smacv2}. \stub{Search for studies that compare multi-agent
results across vehicle-model fidelity before claiming the question is
understudied.}

\section{System}
\label{sec:system}

Ripple began as a Rust port of SCRIMMAGE \cite{demarco2018}; its agreement
with the C++ reference is documented separately \cite{baselinereport}. This
section describes the properties the fidelity question relies on.

\paragraph{One deterministic tick}
Each tick runs fixed phases: generation; autonomy; controllers then motion in
each substep; sensors; interactions; networks; delivery and estimation
(localization then tracking); metrics; and removal, output, and time advance.
Every phase completes before the next begins,
even when entities run on different CPU workers, and output is byte-identical
at 1, 2, and 8 workers and with recording on or off in the core comparisons
\cite{baselinereport}. Sensor random streams use stable entity and plugin
identities; legacy spawn draws retain their reference order. External
autopilots require their own repeatability checks.

\paragraph{Truth, measurements, and belief}
Sensors read physical truth and produce measurements; estimators turn
measurements and received messages into each agent's belief; decisions and
controllers read only belief. A model of higher fidelity therefore reaches a
decision only through what the agent could have measured or been told.

\paragraph{Scale}
Finding nearby agents and routing messages each cost the square of the team
size in the design we started from. Answering spatial queries from a grid, and
routing each message only to its subscribers, made a 4{,}096-agent step 3.8 and
40 times faster with byte-identical output \cite{baselinereport}. That lets a
few expensive vehicles share a world with many cheap ones.

\paragraph{World resources and autopilots}
Terrain is a shared, read-only world resource that motion models query, for
example for landing-gear contact; a sea surface is planned as another
(Section~\ref{sec:platforms}). ArduPilot's software-in-the-loop (SITL)
autopilot flies a vehicle through a lockstep bridge: each simulation step waits
for the autopilot's outputs and answers with the vehicle's state, and several
autopilots run side by side, one per vehicle.

\section{Evidence at Each Integration Stage}
\label{sec:verification}

The C++ comparison establishes a starting point, not a blanket guarantee for
new models. Each integration needs its own external baseline. We distinguish
three questions:
\begin{enumerate}
\item \emph{Did we preserve the source behavior?} Compare a port against its
source on matched inputs and initial conditions. For a bridge, check the
exchange of states and commands, coordinate conversions, and simulated timing
against the external system's interface and a reference run.
\item \emph{Does the model represent the vehicle in the claimed regime?}
Compare against independent measurements where available. Agreement with
another simulator establishes implementation agreement; it does not establish
physical accuracy. Keep calibration data separate from validation data.
\item \emph{Does the added capability matter for the experiment?} Compare
outcomes, failures, and cost under controlled changes in model effects or
control stack. Report the operating regimes actually visited.
\end{enumerate}

For each comparison, retain the source version, model parameters, frames and
units, input sequence, initial state, timestep, and error measures. State the
comparison tolerance and its rationale before interpreting the result. Start
with forces, moments, derivatives, or command outputs where those are the
integration boundary, then compare trajectories and a closed-loop mission.
Preserve failed comparisons and explain intended departures from the source.
An effect disabled in existing missions must preserve their outputs; a
deliberate defect correction needs a documented new baseline.

Table~\ref{tab:integration} ties this method to the implementation roadmap.
These are completion criteria, not claims that the comparisons have run.
Source comparisons belong with each implemented stage, even when physical
validation must wait for data. Each stage also measures execution cost at its
intended team size, so the expense of a few detailed vehicles or external
autopilots is visible.

\begin{table*}[t]
\centering
\small
\caption{External baselines for the integration roadmap. Except for the core
comparison, these are required or planned checks, not completed results.}
\label{tab:integration}
\begin{tabularx}{\textwidth}{@{}P{0.17\textwidth}P{0.37\textwidth}X@{}}
\toprule
Stage & External baseline and comparison & What the evidence can establish \\
\midrule
Simulation core & C++ SCRIMMAGE: matched missions, internal traces, message delivery, and planted defects \cite{baselinereport} & Agreement for tested configurations; a foundation for extensions \\
Fixed-wing tables and autopilot & JSBSim table evaluations on matched inputs; ArduPlane state/actuator exchange and reference mission & Table and bridge behavior; repeatability and closed-loop operation, without claiming Rascal flight accuracy \\
Multirotor dynamics & RotorS motor/rotor outputs sample by sample; Crazyflie parameters from its measured identification & Kernel agreement and sourced parameters; prediction error against flight records is an extension point \\
VTOL & Matched mission against ArduPilot quadplane SITL & Autopilot compatibility; transition-regime accuracy needs measured data such as LA-8's (an extension point) \\
Vessels and waves & MSS Otter sample by sample; hull-point buoyancy in waves as in Gazebo and VRX & Source-model agreement; seakeeping with computed hydrodynamic data is an extension point, distinct from tank or field validation \\
Imported policy (conditional) & Original policy runtime on identical observations; matched normalization, action semantics, and step/reset behavior & Policy and interface agreement before cross-fidelity evaluation \\
\bottomrule
\end{tabularx}
\end{table*}

\section{Fidelity as an Experimental Variable}
\label{sec:fidelity}

To ask whether fidelity changes a result, fidelity must be a controlled
variable: changed alone, with everything else held fixed.

\paragraph{Attributable differences}
Run the same mission, seed, and strategy with a specified model change.
Repeatability makes differences easier to diagnose; attribution also requires
holding other inputs and interfaces fixed. Record controller tuning, decision
periods, sensor and network settings, and initial conditions alongside the
model variant. Core output is byte-identical in the comparisons described in
Section~\ref{sec:system}. Three runs of the runway mission with ArduPlane
(Section~\ref{sec:stall}) also wrote byte-identical output on one machine.
This is evidence for that configuration, not a guarantee for all external
stacks. Across machines and ArduPilot versions it is untested.

\paragraph{Selectable effects}
Each physical effect is an optional parameter of its model, off by default. A
mission that does not set it runs exactly as before, so earlier results and
the verification of the base system \cite{baselinereport} stay valid as
features are added. Where an inherited model is wrong rather than incomplete,
it is fixed and the fix recorded. Vehicles of different
fidelity share one world and the same sensing, network, and estimation
pipeline, so a few expensive vehicles can run among many cheap ones.

\paragraph{Model detail, physical evidence, and control stack}
Table~\ref{tab:ladder} separates choices that a single fidelity score would
conflate. Kinematics and rigid-body dynamics change the model structure;
measured parameters change its evidential basis; an external autopilot changes
the control stack. An autopilot can fly an inaccurate plant. A simple model
with measured parameters can be better supported in its operating regime than
a more detailed model with borrowed coefficients. Comparisons should name the
axis changed and the evidence supporting it.

\begin{table}[t]
\centering
\small
\caption{Independent choices for fidelity experiments in Ripple.}
\label{tab:ladder}
\begin{tabularx}{\linewidth}{@{}P{0.22\linewidth}>{\raggedright\arraybackslash}X>{\raggedright\arraybackslash}X@{}}
\toprule
Axis & Choices & Status \\
\midrule
Motion model & Point agents, constrained kinematics, rigid-body dynamics & Built; multirotor defects under review \\
Physical effects & Stall, contact, motor lag, rotor drag, waves & Stall and terrain contact built; others planned \\
Physical evidence & Borrowed, estimated, or measured parameters; independent validation & Provenance audit; flight-data validation planned \\
Control stack & Native controllers or external autopilot & ArduPlane SITL built; other vehicle stacks planned \\
\bottomrule
\end{tabularx}
\end{table}

\paragraph{The experiment}
For one coordination workload (Section~\ref{sec:workloads}), compare selected
model variants across paired seeds and task schedules, then remove individual
effects within a fixed model and control stack. Keep the strategy's command
interface fixed where possible; report changes in controllers or tuning as a
separate factor. Measure strategy rankings, outcome differences, failure rates,
and computation cost, with uncertainty across runs. Define what outcome
difference would matter for the task before claiming a simpler model suffices.
Record the angles of attack, actuator saturation, or sea conditions visited.
An unchanged ranking alone does not establish equal outcomes or physical
accuracy; a trial outside an effect's operating regime cannot establish that
the effect is generally unnecessary.

\section{Where Model Data Comes From}
\label{sec:provenance}

A model has a structure, its equations, and numbers, its coefficients. Higher
fidelity means something only if both come from stated sources. We checked the
sources of the models we use or plan to use, in September and October 2026
(Table~\ref{tab:provenance}).

\begin{table*}[t]
\centering
\small
\caption{Where the structure and the numbers of open vehicle models come from,
as checked in September and October 2026.}
\label{tab:provenance}
\begin{tabularx}{\textwidth}{@{}P{0.17\textwidth}P{0.17\textwidth}P{0.2\textwidth}X@{}}
\toprule
Source & Vehicle & Structure & Numbers \\
\midrule
JSBSim Rascal, the source of our fixed-wing coefficients & 6~kg radio-controlled aircraft & Coefficient tables against angle of attack & Adapted from FlightGear for ArduPilot testing; its README reads ``we are still at pre-pre-pre-alpha with this model'', and its author field ``Author Name'' \\
ArduPilot SITL plane \cite{ardupilotsitl} & Generic & Linear lift blended into a flat plate past a stall angle & Its source describes it as ``Not aerodynamically accurate, just enough to be able to debug control logic'' \\
RotorS \cite{furrer2016rotors} & AscTec Hummingbird quadrotor; Firefly hexacopter & Motor lag, rotor drag, rolling moment & Hummingbird mass, arm, and inertia from a published paper; all seven rotor and motor coefficients are identical, digit for digit, for both vehicles, and the rolling-moment coefficient is $10^{-6}$ \\
PX4 Gazebo motor model \cite{px4gazebo} & Generic & RotorS's, plus thrust falling linearly with axial airspeed to zero at 25~m/s & The thrust reduction is commented ``XXX this has to be modelled better'' \\
MSS Otter \cite{fossen2021handbook} & 2~m uncrewed catamaran & Six-degree-of-freedom manoeuvring and hydrostatics & Dimensions and masses of the real vehicle; added mass as fixed fractions of mass (0.1, 1.5, and 1.0 in surge, sway, and heave); sway and yaw time constants of 1~s \\
Crazyflie identification \cite{forster2015crazyflie} & 27~g quadrotor & Thrust and torque against rotor speed; inertia; motor dynamics & Measured on the vehicle with a force/torque sensor and a laser tachometer; 172 public flights with motion capture \cite{ullah2026nanobench} \\
NeuroBEM \cite{bauersfeld2021neurobem} & 752~g racing quadrotor & Blade-element rotor model with a learned residual & 96 flights, 1~h~15~min of motion capture and motor speeds, up to 18~m/s \\
NASA LA-8 \cite{simmons2021la8,simmons2023dissertation,north2021la8,simmons2021propeller,geuther2020la8} & 28~kg tandem tilt-wing, eight propellers & Aero-propulsive response surfaces at eight airspeeds & Measured mass properties and a propeller model from 0 to 180$^\circ$ incidence; whole-aircraft coefficients printed at one of the eight airspeeds (16.5~m/s); about 300 pages of wind-tunnel plots \\
\bottomrule
\end{tabularx}
\end{table*}

Three observations follow. Structure is easy to obtain: every source publishes
its equations. Numbers measured on the vehicle being simulated are rare, and
they come from system identification of one airframe. Borrowed numbers are hard
to see: nothing in RotorS's vehicle files marks its rotor coefficients as
shared, and nothing in the fixed-wing model we inherited records that its
coefficients came from a model its authors called pre-alpha. For boats in
waves, hydrodynamic coefficients are usually computed by panel-method solvers
such as Capytaine \cite{ancellin2019capytaine} rather than measured. Each model
in Ripple should therefore state the source of its structure and of its
numbers, and the regime in which they were measured, and claim no more fidelity
than those sources support.

\stub{Record each model's sources in the book's model reference.}

\section{A First Extension: Stall}
\label{sec:stall}

FixedWing6DOF, the fixed-wing model we started from, computes lift as $C_{L0} +
C_{L\alpha}\alpha$ plus
rate and elevator terms at any angle of attack $\alpha$, so the wing never
stalls. Its only drag that grows with angle, for the Rascal, is induced drag,
$C_L^2/(\pi A e)$, which follows the lift: a lift curve that falls past the
stall would also lose drag there, the opposite of a stalled wing.

\paragraph{What we added}
Two optional tables against angle of attack, after JSBSim. A lift table
replaces $C_{L0} + C_{L\alpha}\alpha$; a drag table replaces the zero-lift drag
$C_{D0} + C_{D\alpha}|\alpha|$, and induced and elevator drag are still added.
Values are interpolated linearly and held at the ends, as JSBSim's tables are.
JSBSim's Rascal lift table equals the mission's linear lift up to a peak at
0.23~rad (13.2$^\circ$), which puts the level stall near 8.3~m/s; its drag
table rises from 0.028 at zero to 1.5 at $\pm90^\circ$. Without tables the
model is the C++ one, and its two comparison missions still agree with C++ to
rounding \cite{baselinereport}. Spin and drag from sideslip are not modeled; a
deep stall or large sideslip counts as a failure already.

\paragraph{Does it change a landing?}
ArduPlane SITL flies the Rascal off a runway on surveyed terrain, around a
circuit, and lands it with a 14$^\circ$ flare. With linear lift, the landing
bounced. We flew it three ways: as committed, with the lift table, and with
both tables (Table~\ref{tab:stall}; \texttt{paper/data/stall-landing}). A
repeat of the unmodified mission was byte-identical to the first run, so the
differences between variants are effects of the tables, not run-to-run noise.

\IfFileExists{generated/stall.tex}{% Generated by paper/data/make_tables.py; do not edit.
\begin{table}[t]
\centering
\small
\caption{ArduPlane's runway landing of the Rascal with and without
stall tables: speed and sink rate at first touchdown, highest bounce, widest
offset from the centreline after touchdown, and largest airborne angle of
attack from 15~s before touchdown.}
\label{tab:stall}
\begin{tabular}{@{}lrrrrr@{}}
\toprule
Lift and drag & m/s & sink & bounce & offset & $\alpha$ \\
\midrule
Linear & 15.36 & 1.46 & 2.03 & 1.56 & 11.2$^\circ$ \\
Lift table & 15.36 & 1.45 & 2.04 & 1.71 & 10.8$^\circ$ \\
Both tables & 15.09 & 1.53 & 1.88 & 1.23 & 8.9$^\circ$ \\
\bottomrule
\end{tabular}
\end{table}
}{}

The lift table changed the landing by centimetres. Airborne, the angle of
attack never passed 11.2$^\circ$, below the 13.2$^\circ$ stall; only rolling
tail-down on the wheels, above 5~m/s, did it reach 14.5$^\circ$. The drag table
changed more, because it adds drag at every angle, not only past the stall:
touchdown 0.27~m/s slower, a bounce 0.15~m lower, and a rollout within 1.23~m
of the centreline instead of 1.56~m. The bounce itself remains. Touchdown at
about 15~m/s, nearly twice the stall speed, and main wheels ahead of the centre
of gravity, which pitch the nose up on contact, are likely causes; neither is
tested.

\paragraph{What it shows}
This mission exercises the integrated tables and autopilot, but it does not
establish the accuracy of post-stall flight or a multi-agent consequence of
stall. The recorded airborne landing interval stays below the stall angle,
and the larger change comes from drag at small angles. Direct comparison with
the source table evaluations remains a separate check
(Section~\ref{sec:verification}). The lesson is to measure the regime and
state exactly which claim a trial supports.

\section{Planned Platforms}
\label{sec:platforms}

Each planned model reaches the level of an established open simulator
(ArduPilot and PX4 SITL, RotorS, Fossen's MSS): structure from a documented
source, every number's origin stated, and a check against the source
implementation. Going further needs measured data for one real vehicle, which
we do not collect. Instead each model is built so that such data plugs in: a
header on every motion model states its source, what it was validated
against, its valid range, what it leaves out, and what data would take it
further; effects left out are named zero terms in the force sums; and geometry
and coefficients are parameters (rotor axes as vectors, aerodynamics as tables),
so a measured model arrives as data or one function.

\paragraph{Multirotors}
Reviewing the inherited multirotor model for this work found three defects,
present in the simulator it was ported from as well. Position
is integrated from a body-frame velocity that starts at zero, while the
reported world velocity is integrated separately, so the two can disagree.
Moment arms use the full three-dimensional rotor offset, so a rotor's height
above the centre of gravity adds leverage, and configured rotor tilt is
ignored. Thrust drops to zero only when a rotor's speed equals the minimum
exactly, while its yaw torque remains. These will be fixed first. Motor lag,
rotor drag, and the
rolling moment will follow RotorS \cite{furrer2016rotors} as optional effects.
Rotor drag, linear in airspeed, dominates at typical multirotor speeds
\cite{martin2010accelerometer}; the inherited quadratic body drag remains as a
second term for high speed. RotorS's own motor constants carry no source, so
the real-vehicle parameter set comes from the Crazyflie's measured
identification \cite{forster2015crazyflie}, each value cited, beside the
inherited defaults kept for the C++ comparison. Replaying public flights with
recorded motor commands \cite{ullah2026nanobench} would give a stated
prediction error; that validation is an extension point. A controller comes
first: no controller yet flies the model closed loop.

\paragraph{VTOL aircraft}
Transitions between hover and wingborne flight pass through the angles at which
a wing stalls, so they exercise the tables of Section~\ref{sec:stall}. A
quadplane, FixedWing6DOF with lift rotors flown by ArduPilot's quadplane SITL,
would reach the fidelity of ArduPilot's and PX4's own simulators, neither of
which models propeller wash over the wing \cite{ardupilotsitl,px4gazebo}.
Rotor axes are fixed vectors, so a tilt-rotor needs only a control input for
them, and wing slipstream is a named zero term. NASA's LA-8 is the only open
source of measured transition data we found; its whole-aircraft model is
published at one of eight airspeeds
\cite{simmons2021la8,simmons2023dissertation},
so a measured transition model depends on the rest being released.

\paragraph{Surface vessels}
There is no surface vessel yet. Fossen's
Otter \cite{fossen2021handbook} gives calm-water manoeuvring and hydrostatic
restoring forces. Waves fit the architecture as a world resource beside
terrain: a sum of spectral components whose phases are drawn once from the
mission seed is a pure function of position and time, so every worker reads the
same sea and output stays byte-identical. Wave loads are buoyancy at points
on the hull, the level of Gazebo's and VRX's boats \cite{bingham2019vrx}. Motion
transfer functions and linear seakeeping with radiation memory, using
coefficients computed by a panel solver \cite{ancellin2019capytaine}, are the
documented next steps.

\stub{For each implemented platform, report the external comparison, the
remaining physical uncertainty, and one representative closed-loop mission.
Use selected platforms for the coordination study; every platform need not
support every workload before a bounded contribution can be reported.}

\section{Heterogeneous Fleets}
\label{sec:fleets}

The platforms matter most together. An uncrewed boat carries a multirotor that
surveys ahead and lands back on its deck in waves; fixed-wing aircraft cover a
wide area while multirotors inspect what they find; all share what they learn
over links that drop messages. Each combination asks the fidelity question
across domains: does a deck landing that succeeds in calm water fail in a
modeled sea, and does that change which coordination strategy wins?

Bringing the platforms into one world needs:
\begin{itemize}
\item \emph{One tick across domains.} Vessel, multirotor, and aircraft dynamics
advance in the same phases and substeps, so their interactions are ordered
deterministically.
\item \emph{Shared world resources.} Terrain for aircraft and wheels, the sea
surface for vessels, and wind, each a read-only function that every model and
sensor queries.
\item \emph{Contact between vehicles.} Landing gear touches terrain today; a
multirotor on a moving deck needs contact between two vehicles' dynamics.
\item \emph{An autopilot per vehicle.} ArduPlane, ArduCopter, and ArduRover in
lockstep, one per vehicle; two ArduPlane instances already run side by side.
\item \emph{One information pipeline.} The same sensing, network, and
estimation stages for every domain, so a boat's report reaches an aircraft's
belief with the same delay and loss rules.
\item \emph{Mixed fidelity.} A few high-fidelity vehicles among many simple
ones, chosen per vehicle in the mission.
\end{itemize}

\stub{Begin with independently moving vehicles sharing information in one
world. Moving-deck contact is a further integration with its own mechanical
baseline, not a prerequisite for the current roadmap.}

\section{Coordination Workloads}
\label{sec:workloads}

The integrations need representative multi-agent use, followed by controlled
fidelity comparisons. Information sharing is the first workload. Planning
computation and team scale are later options, not additional prerequisites for
completing the vehicle roadmap or reporting an integration.

\subsection{Sharing Information}
\label{sec:rq1}
\label{sec:rq4}

When does sharing information improve team performance? Kimera-Multi and RACER
motivate asking what robots should share and when its value outweighs its cost
\cite{tian2022kimera,zhou2023racer}. The estimation stage
(Section~\ref{sec:system}) makes the question testable: decisions read only
belief, so an experiment controls exactly what each robot knows. Compare local
decisions, periodic sharing, and selective sharing on one exploration or
target-visiting task under a fixed message budget, measuring coverage,
duplicated work, and stale decisions, with ideal information as a separate
reference. Fidelity enters through how long a vehicle takes to act on what it
learns: a quadrotor with motor lag or a boat in waves reaches a reported target
later than a point agent, which may change when sharing pays off.

\subsection{Planning Computation}
\label{sec:planning}

When does more planning computation improve the outcome? Accelerated
simulators make evaluating many alternatives cheap
\cite{makoviychuk2021isaacgym,bettini2022vmas,mittal2025isaaclab}; the question
is whether doing so changes the decision that can be made in time. A planner
plugin can already score its own candidates. Batches across entities, and a
rule for when a slow decision takes effect, would need an optional shared
interface, with prescribed delays so that GPU completion order never determines
simulated behavior. Fidelity enters through the planner's model: candidates
scored with a point model and flown by a vehicle that lags or stalls may rank
differently from candidates scored with the vehicle's own model.

\subsection{Team Scale}
\label{sec:rq3}

When does a larger team stop completing more work? RHCR and RACER motivate
measuring sustained throughput as demand and team size grow
\cite{li2021rhcr,zhou2023racer}. Entity generation and removal exist, but task
arrival, ownership, completion, and reassignment do not; a reusable task model
with saved arrival schedules would begin as plugins. Compare allocation methods
under identical arrivals while increasing demand, team size, and interruptions,
measuring completed tasks per unit time, waiting time, and backlog. Fidelity
enters through crowding: keeping separation costs a point agent nothing and a
multirotor time.

\stub{Pilot: every message delivery is already traced, so sharing information
is the natural first workload.}

\section{Learned Policies}
\label{sec:learning}

Learning is worth its cost for decisions that classical strategies cannot
hand-code: weighing many tracks and their uncertainty, teammates' messages,
vehicle limits, and timing at once. It is not needed to re-solve what a
classical strategy already does well. A policy for such a decision can learn
only from the inputs it will face, so it must train in worlds with the same
sensing, networks, belief, and vehicles. Those are Ripple's own worlds, and
learned policies belong in the fidelity question with particular force,
because a policy can exploit a model's simplifications.

\paragraph{Training on Ripple's worlds}
Ripple contains no training code; an outside learner drives its worlds through
a standard multi-agent interface
\cite{terry2021pettingzoo,bettini2024benchmarl}.
The policy makes a high-level decision (a target, a waypoint, a message, a
hand-off) at about 1~Hz, and the classical autonomy, controllers, and motion
carry it out. Learning decisions over existing behaviors, rather than
low-level control, cuts the samples needed by orders of magnitude and keeps the
vehicle within its limits. During training a critic may read truth while each
agent's policy sees only its own belief and messages, a separation Ripple
already enforces. Rewards come from the metrics plugins, so the training reward
and the evaluation score agree, and one definition of observations and actions
serves both training and evaluation. Throughput, decisions per second per core
at each fidelity level, is unmeasured and is the first thing to measure.

\paragraph{CPU worlds and GPU simulators}
Simulating on CPUs while a GPU trains the policy is the established
arrangement for complex environments, and current libraries are built to make
it fast \cite{petrenko2020samplefactory,suarez2025pufferlib}. GPU-resident
simulators run thousands of worlds on one device at millions of steps per
second
\cite{makoviychuk2021isaacgym,mittal2025isaaclab,rutherford2023jaxmarl,shacklett2023madrona},
but they reach that speed with batchable models, and most still simplify the
sensing, belief, and communication these decisions depend on, even as some add
vehicle fidelity (Section~\ref{sec:related}). The
question is therefore not CPU against GPU but whether the cheaper world
produces a policy that still works in the richer one. We plan a hybrid: train
cheaply wherever possible, including in a GPU simulator, then fine-tune and
evaluate in Ripple against classical strategies as fidelity rises
(Section~\ref{sec:fidelity}). A policy trained elsewhere enters as a frozen
policy behind an autonomy plugin; its observation normalization, frames, action
meanings, decision period, and termination rules must match its training task.
If a cheaply trained policy holds up at high fidelity, that is a result too.
Transfer across fidelity within one simulator complements BVR Sim's transfer
across aircraft within one flight model \cite{sun2026bvrsim}. One standard
benchmark task reproduced in Ripple would anchor its learning results to
published ones. Table~\ref{tab:learning} lists the options; none is built, and
each is independent of the others.

\begin{table*}[t]
\centering
\small
\caption{Options for learning and GPU computation. Each is independent of the
others.}
\label{tab:learning}
\begin{tabularx}{\textwidth}{@{}P{0.25\textwidth}>{\raggedright\arraybackslash}X>{\raggedright\arraybackslash}X@{}}
\toprule
Option & What it provides & Fit \\
\midrule
CPU simulation, imported CPU policy & Learned-versus-classical evaluation across fidelity & Smallest useful integration \\
CPU worlds, external GPU learner & Training on this testbed's own information, task, and vehicle rules & Decisions that need the full information; hierarchical over classical behaviors \\
GPU simulator, then Ripple & Cheap training, then fine-tuning and evaluation at higher fidelity & Tests whether the cheap world was enough \\
CPU simulation, batched GPU inference & Faster evaluation of large policies or many agents & When inference dominates the step \\
GPU candidate trajectories in one planner & More planning alternatives per decision & Inside one optional plugin \\
GPU physics or sensing kernels & Acceleration of a measured expensive phase & When the gain covers packing, transfers, and synchronization \\
Fully GPU-resident worlds & Large rollout batches & A new backend with restricted models and its own verification \\
\bottomrule
\end{tabularx}
\end{table*}

Any GPU computation must keep the determinism of Section~\ref{sec:system}:
repeated evaluation on a pinned backend and device, invariant to worker count.
Assembling each batch in entity order removes batch composition as a source of
variation, which dominates nondeterminism in LLM inference
\cite{he2025nondeterminism}; the kernels and any sampling must also be
deterministic.

\stub{Measure decision throughput in the information-sharing workload at
several fidelity levels; choose the decision to learn and record why classical
strategies are not enough for it; then report learning curves and the
fidelity-transfer comparison.}

\section{Discussion}
\label{sec:discussion}

\subsection{What Makes an Integration Useful?}

An integration is useful when it enables an intended experiment and has enough
evidence to interpret its results. That evidence includes agreement with an
external baseline, clear physical limits, and measured execution cost. It
does not require every new effect to change a strategy ranking. Conversely,
an impressive trajectory or a changed ranking cannot substitute for checking
that the imported component still behaves as intended.

The roadmap can therefore advance stage by stage. Each stage contributes a
working capability and a bounded comparison. Studies across fidelity levels
then establish where the capability affects conclusions and where a cheaper
model is sufficient within a stated margin. Coordination studies still need
held-out layouts, task schedules, or team sizes, and removal of communication
or feedback to check that it matters \cite{ellis2023smacv2}.

\subsection{Limits}

No model is yet checked against a vehicle's measured flight data; the stall
tables carry JSBSim's Rascal values, which their own README calls pre-alpha.
The stall result is one mission, one aircraft, and one autopilot
configuration, on one machine; its byte-repeatability with ArduPilot is shown
there and not across machines or ArduPilot versions. Effects that are not
modeled, such as spin, propeller wash over a wing, and waves between vessels,
bound what the experiments can test. Agreement of the base system with C++
SCRIMMAGE is established for the models and configurations tested there
\cite{baselinereport}; extensions require their own comparisons. No
multi-agent fidelity result is reported yet.

\section{Conclusion}

\stub{State the completed integrations and the external evidence supporting
each. Then report what the selected fidelity experiments establish, including
unchanged outcomes and cost. Separate source agreement from physical accuracy,
and end with the next integration or measurement justified by the results.}

\section*{Generative AI Disclosure}

Claude (Anthropic) assisted with implementation, running the experiments,
checking sources for the provenance audit, and drafting text from retained
data. OpenAI Codex assisted with revising the paper's framing and integration
evaluation plan. \stub{Update this statement to reflect actual assistance and
author
review, following the selected venue's policy.}

\bibliographystyle{plainnat}
\bibliography{../references}

\end{document}
